\usepackage[english]{babel}
\usepackage[T1]{fontenc}
\usepackage[utf8]{inputenc}
\usepackage{microtype}
\usepackage{stmaryrd}
\usepackage[dvipsnames]{xcolor}% needs to be loaded before todonotes, since todonotes loads xcolor without options.
\usepackage{subfiles}

\usepackage[textsize = footnotesize, colorinlistoftodos]{todonotes}
\newcommand{\todolater}[1]{\todo[color=gray!30,linecolor=lightgray,size=scriptsize,textcolor=gray,bordercolor=lightgray]{\textbf{Later:} #1}} %Notes that can be ignored for now
\newcommand{\fnote}[1]{\todo[color=yellow!50,linecolor=lightgray,size=scriptsize]{\textbf{FA:} #1}}
\newcommand{\fnoteil}[1]{\todo[inline,color=yellow!50,linecolor=lightgray,size=large]{\textbf{FA:} #1}}
\newcommand{\tnote}[1]{\todo[color=pink!50,linecolor=lightgray,size=scriptsize]{\textbf{TP:} #1}}
% \newcommand{\tnoteil}[1]{\todo[inline,color=pink!50,linecolor=lightgray,size=large]{\textbf{Tommy:} #1}}
\newcommand{\tnoteil}[1]{

{\color{red} \textbf{TP:} #1}

}
\newcommand{\todoin}[2][]{\todo[inline,#1]{#2}}
\newcommand{\comm}[1]{{\color{blue} [ #1 ]}}


% \usepackage{showlabels}
% \renewcommand{\showlabelfont}{\small\ttfamily\color{orange}}



%font
\usepackage{eucal}
\usepackage[shortlabels]{enumitem}
\usepackage[super]{nth}


% \usepackage[backend=biber,
% style=ext-alphabetic,
% sorting=nyt,% sort by name, year, title
% giveninits=true,% render all first names as initials.
% isbn=false,
% % url=false,
% maxalphanames=4,
% % minalphanames=3,
% date=year,
% maxbibnames=99,
% backref=false,
% articlein=false]{biblatex}
% \addbibresource{references.bib}
% \usepackage{csquotes}

% to remove the "In:" in articles
%\renewbibmacro{in:}{%
%	\ifentrytype{article}{}{\printtext{\bibstring{in}\intitlepunct}}}

% amspackages
\usepackage{amsmath}
\usepackage{amssymb}
\let\openbox\relax
\usepackage{amsthm}
\usepackage{amsfonts}
\usepackage{graphicx}

%mathpacks
\usepackage{dsfont} %blackboard bold 1
\usepackage{mathtools}
\usepackage{xparse}
\usepackage{tikz-cd}
\usepackage{quiver}
\usepackage{nccmath}
%\tikzset{%adjunction symbol
%	symbol/.style={%
%		draw=none,
%		every to/.append style={%
%			edge node={node [sloped, allow upside down, auto=false]{$#1$}}}
%	}
%}

% tikz
\usepackage{tikz}
\usepackage{tabularx}
\newcolumntype{Y}{>{\centering\arraybackslash}X}
% \usetikzlibrary{shapes.geometric, arrows}
\newcommand{\nwidth}{4cm}
\tikzstyle{no} = [rectangle, rounded corners, 
minimum width=\nwidth, 
minimum height=1cm,
text width = \nwidth,
align=center,
draw=black]

\tikzstyle{arrowv} = [implies-,double equal sign distance,shorten >=0.15cm,shorten <=0.15cm]
\tikzstyle{arrow} = [implies-,double equal sign distance,shorten >=0.05cm,shorten <=0.05cm]

%hyperref?
\usepackage[pdfborder={0 0 0}]{hyperref}
\hypersetup{
	colorlinks=true,
	allcolors=.,
	bookmarksdepth = 3
%	linkcolor={BrickRed},
%	citecolor={PineGreen},
%	urlcolor={black},
}
\setcounter{tocdepth}{2}
\makeatletter
\def\@tocline#1#2#3#4#5#6#7{\relax
  \ifnum #1>\c@tocdepth
  \else
    \par \addpenalty\@secpenalty\addvspace{#2}%
    \begingroup \hyphenpenalty\@M
    \@ifempty{#4}{%
      \@tempdima\csname r@tocindent\number#1\endcsname\relax
    }{%
      \@tempdima#4\relax
    }%
    \parindent\z@ \leftskip#3\relax \advance\leftskip\@tempdima\relax
    \rightskip\@pnumwidth plus4em \parfillskip-\@pnumwidth
    #5\leavevmode\hskip-\@tempdima #6\nobreak\relax
    \leaders\hbox{$\m@th\mkern4.5mu\hbox{.}\mkern4.5mu$}\hfil
    \hbox to\@pnumwidth{\@tocpagenum{#7}}\par
    \nobreak
    \endgroup
  \fi}
\renewcommand\l@subsection{\@tocline{2}{0pt}{2.5pc}{5pc}{}}
\makeatother
\usepackage[capitalize,nameinlink,nosort]{cleveref}
\crefformat{equation}{#2(#1)#3}
\Crefformat{equation}{#2(#1)#3}
\crefname{enumi}{}{}
\crefname{construction}{Construction}{Constructions}

%---------------------------------------------------------
% Dit blok fixt een probleem met de counters wat ervoor zorgde that cleveref niet meer werkte. Lijkt nu allemaal weer goed te gaan
\usepackage{aliascnt}
\let\oldtheorem\newtheorem
\RenewDocumentCommand{\newtheorem}{s m o m O{}}{%
\IfBooleanTF{#1}%
{\oldtheorem{#2}{#4}}%
{\IfNoValueTF{#3}{\oldtheorem{#2}{#4}[#5]}%
{\newaliascnt{#2}{#3}%
\oldtheorem{#2}[#2]{#4}%
\aliascntresetthe{#2}}}}
%---------------------------------------------------------

% theorems etc.
\theoremstyle{plain}
\newtheorem{theorem}{Theorem}[section]
\newtheorem{lemma}[theorem]{Lemma}
\newtheorem{proposition}[theorem]{Proposition}
\newtheorem{corollary}[theorem]{Corollary}
\newtheorem{fact}[theorem]{Fact}
\newtheorem*{theorem*}{Theorem}
\newtheorem{thmint}{Theorem}
\newtheorem{corint}[thmint]{Corollary}
\renewcommand*{\thethmint}{\Alph{thmint}}
\renewcommand*{\thecorint}{\Alph{thmint}}

\theoremstyle{definition}
\newtheorem{definition}[theorem]{Definition}
\newtheorem{remark}[theorem]{Remark}
\newtheorem{example}[theorem]{Example}
\newtheorem{notation}[theorem]{Notation}
\newtheorem*{notation*}{Notation}
\newtheorem{construction}[theorem]{Construction}
\newtheorem{warning}[theorem]{Warning}
\newtheorem{observation}[theorem]{Observation}

%---(nicer underlines)-------------------------------------------------------------------START
\makeatletter
\usepackage[normalem]{ulem}
\usepackage{contour}

\renewcommand{\ULdepth}{1.8pt} % how low to make the underline
\contourlength{0.5pt} % how big to make the gap in words like map

\makeatletter
% usage: there are two ways to use this 
% \myuline{text} puts an underline under text
% \myuline[7em]{text} puts an underline under text but shortens it by 7em; of course 7em can be changed. This is useful for stuff like \({\myuline[0.5pt]{\mathrm{Map}}}_G\), vs \({\myuline{\mathrm{Map}}}_G\)
\DeclareRobustCommand{\myuline}[2][0pt]{%
  \ifmmode
    % draw a slightly shorter underline
    \uline{\hphantom{#2}\kern-#1}%
    % restore overall width so following text lines up
    \kern#1%
    % overlay the actual content, with proper math style
    \mathllap{\mathpalette\my@cont@{#2}}%
  \else
    \uline{\phantom{#2}\kern-#1}%
    \kern#1%
    \llap{\contour{background}{#2}}%
  \fi
}
% #1 = current math style, #2 = content
\newcommand{\my@cont@}[2]{\contour{background}{\mbox{$\m@th#1#2$}}}
\makeatother
%------------------------------------------------------------------------------------------END

\usepackage{relsize}
\usepackage[bbgreekl]{mathbbol}
\usepackage{amsfonts}
\DeclareSymbolFontAlphabet{\mathbb}{AMSb} %to ensure that the meaning of \mathbb does not change
\DeclareSymbolFontAlphabet{\mathbbl}{bbold}





%% COMMANDS AND OPERATORS

%% operators
\DeclareMathOperator{\im}{Im}
\DeclareMathOperator{\coker}{coker}
\DeclareMathOperator{\Hom}{Hom}
\DeclareMathOperator{\End}{End}
\DeclareMathOperator{\Aut}{Aut}
\newcommand{\id}{\mathrm{id}}
\DeclareMathOperator*{\colim}{colim}
\DeclareMathOperator*{\laxlim}{laxlim}
\DeclareMathOperator*{\oplaxlim}{oplaxlim}
\DeclareMathOperator*{\laxcolim}{laxcolim}
\DeclareMathOperator{\Spec}{Spec}
\DeclareMathOperator{\Ad}{Ad}
\DeclareMathOperator{\ad}{ad}
\DeclareMathOperator{\val}{ord}
\DeclareMathOperator{\ev}{ev}
\DeclareMathOperator{\inv}{inv}
\DeclareMathOperator{\dv}{div}
\DeclareMathOperator{\Gal}{Gal}
\DeclareMathOperator{\GL}{GL}
\DeclareMathOperator{\SL}{SL}
\DeclareMathOperator{\PGL}{PGL}
\DeclareMathOperator{\chr}{char}
\DeclareMathOperator{\Cl}{Cl}
\DeclareMathOperator{\Div}{Div}
\DeclareMathOperator{\Pic}{Pic}
\DeclareMathOperator{\Br}{Br}
\DeclareMathOperator{\Kth}{K}
\DeclareMathOperator{\Tor}{Tor}
\DeclareMathOperator{\Ext}{Ext}
\DeclareMathOperator{\free}{free}
\DeclareMathOperator{\forget}{forget}
\DeclareMathOperator{\cofree}{cofree}
\DeclareMathOperator{\indec}{indec}
\DeclareMathOperator{\Free}{Free}
\newcommand{\BFree}{\mathsf{BFree}}
\newcommand{\oo}{\mathsf{o}}
\DeclareMathOperator{\LMod}{LMod}
\DeclareMathOperator{\BMod}{BMod}
\DeclareMathOperator{\Mod}{Mod}
\DeclareMathOperator{\Alg}{Alg}
\DeclareMathOperator{\coAlg}{coAlg}
\DeclareMathOperator{\CAlg}{CAlg}
\DeclareMathOperator{\coCAlg}{coCAlg}
\DeclareMathOperator{\Map}{Map}
\DeclareMathOperator{\Sym}{Sym}
\DeclareMathOperator{\Shv}{Shv}
\DeclareMathOperator{\Tens}{Tens}
\DeclareMathOperator{\fib}{fib}
\DeclareMathOperator{\cof}{cof}
\DeclareMathOperator{\Ann}{Ann}
\DeclareMathOperator{\Th}{Th}
\DeclareMathOperator{\cross}{cr}
\DeclareMathOperator{\segNerve}{N}
\DeclareMathOperator{\cosk}{cosk}
\DeclareMathOperator{\Equiv}{Equiv}
\DeclareMathOperator{\Bahr}{Bar}
\DeclareMathOperator{\Cobar}{Cobar}
\DeclareMathOperator{\const}{const}
\DeclareMathOperator{\Twr}{Tw^r}
\DeclareMathOperator{\Twl}{Tw^l}
\DeclareMathOperator{\Ar}{Ar}
\DeclareMathOperator{\Lan}{Lan}
\DeclareMathOperator{\Comp}{Comp}
\DeclareMathOperator{\Sq}{Sq}
\DeclareMathOperator{\Ver}{Vert}
\DeclareMathOperator{\Hor}{Hor}
\DeclareMathOperator{\CMon}{CMon}
\newcommand{\Monad}{\mathrm{Monad}}
\newcommand{\whisk}{\cdot}

%comment
\newcommand{\hide}[1]{}

%enumerate upright
\newlist{numberenum}{enumerate}{1}
\setlist[numberenum]{\upshape(\arabic*)}

\newcommand{\Tommy}[2]{\llbracket #1,#2\rrbracket}

%% categories
\newcommand{\Fun}{\mathrm{Fun}}
\newcommand{\twomap}{\myuline[1pt]{\mathrm{Map}}}
\newcommand{\Day}{\mathrm{Fun}}
\newcommand{\Lax}{\mathrm{Lax}}
\newcommand{\Oplax}{\mathrm{Oplax}}
\newcommand{\UnitLax}{\mathrm{Lax}_\mathrm{un}}
\newcommand{\UnitOplax}{\mathrm{Oplax}_\mathrm{un}}
\newcommand{\FunSeg}{\mathrm{Fun}^\mathrm{Seg}}
\newcommand{\inert}{\mathrm{int}}
\newcommand{\Inert}{\mathrm{Inert}}
\newcommand{\activ}{\mathrm{act}}
\newcommand{\elem}{\mathrm{el}}
\newcommand{\idle}{\mathrm{idle}}
\newcommand{\allmaps}{\mathrm{all}}
\newcommand{\lFun}{\mathrm{Fun}^{\mathrm{L}}}
\newcommand{\mlFun}{\mathrm{Fun}^{\mathrm{(L)}}}
\newcommand{\Funtwo}{\mathrm{Fun}_2}
\newcommand{\Funtwofl}{\mathrm{Fun}^\fl_2}
\newcommand{\Fin}{\mathcal{F}\mathrm{in}}
\newcommand{\Surj}{\mathrm{Surj}}
\newcommand{\Assoc}{\mathrm{Assoc}}
\newcommand{\Ab}{\mathcal{A}\mathrm{b}}
\newcommand{\grAb}{\mathrm{gr}\mathcal{A}\mathrm{b}}
\newcommand{\Sp}{\mathrm{Sp}}
\newcommand{\hSp}{\mathrm{hSp}}
\newcommand{\presl}{\mathcal{P}\mathrm{r}^{\mathrm{L}}}
\newcommand{\presr}{\mathcal{P}\mathrm{r}^{\mathrm{R}}}
\newcommand{\Spc}{\mathcal{S}}
\newcommand{\Set}{\mathcal{S}\mathrm{et}}
\newcommand{\pSpc}{\mathcal{S}_{\ast}}
\newcommand{\PreCat}{\mathcal{P}\mathrm{re}\mathcal{C}\mathrm{at}}
\newcommand{\Cat}{\mathcal{C}\mathrm{at}}
\newcommand{\XCat}{\mathcal{X}\text{-}\mathcal{C}\mathrm{at}}
\newcommand{\Catco}[1]{\Cat_{(\infty,1)}^{#1\mhyphen\mathrm{cocomp}}}
\newcommand{\LCat}{\mathcal{C}\mathrm{at}^\mathrm{L}}
\newcommand{\LTop}{\mathcal{L}\mathrm{Top}}
\newcommand{\RTop}{\mathcal{R}\mathrm{Top}}
\newcommand{\Cattwo}{\mathcal{C}\mathrm{at}_2}
\newcommand{\Cattwofl}{\mathcal{C}\mathrm{at}^{\mathrm{fl}}_2}
\newcommand{\Catinfty}{\mathcal{C}\mathrm{at}}
\newcommand{\DblCat}{\mathcal{D}\mathrm{bl}\mathcal{C}\mathrm{at}}
\newcommand{\DblCatlax}{\mathcal{D}\mathrm{bl}\mathcal{C}\mathrm{at}^\mathrm{lax}}
\newcommand{\DblCatulax}{\mathcal{D}\mathrm{bl}\mathcal{C}\mathrm{at}^\mathrm{ulax}}
\newcommand{\coCart}{\Cocart}
\newcommand{\Cocart}{\mathcal{C}\mathrm{ocart}}

\newcommand{\LoccoCart}{\mathcal{L}\mathrm{oc}\mathrm{co}\mathcal{C}\mathrm{art}}
\newcommand{\Cart}{\mathcal{C}\mathrm{art}}
\newcommand{\Fib}{\mathcal{F}\mathrm{ib}}
\newcommand{\PFree}{\mathsf{PFree}}

\newcommand{\Factcat}{\mathcal{F}\mathrm{act}}
\newcommand{\CSS}{\mathcal{CS}\mathrm{eg}}
\newcommand{\twoSeg}{\mathcal{CS}\mathrm{eg}^2}
\newcommand{\MonCat}{\mathcal{M}\mathrm{on}}
\newcommand{\SymMonCat}{\mathcal{S}\mathrm{ym}\mathcal{M}\mathrm{on}}
\newcommand{\Modl}{\widehat{\Mod}}
\newcommand{\Vect}{\mathrm{Vect}}
\newcommand{\Lie}{\mathrm{Lie}}
\newcommand{\Ho}{\mathrm{Ho}}
\newcommand{\Ch}{\mathrm{Ch}}
\DeclareMathOperator{\Env}{Env}
\DeclareMathOperator{\UnitEnv}{Env^\mathrm{un}}
\DeclareMathOperator{\Span}{Span}
\DeclareMathOperator{\BigSpan}{\overline{Span}}
\newcommand{\coEnv}{\mathrm{coEnv}}
\newcommand{\LaxMor}{\overline{\mathbb{M}\mathrm{or}}}
\newcommand{\Mor}{\mathbb{M}\mathrm{or}}
\newcommand{\op}{\mathrm{op}}
\newcommand{\coop}{\mathrm{coop}}
\newcommand{\co}{\mathrm{co}}

\newcommand{\rev}{\mathrm{rev}}
\newcommand{\Seq}{\mathrm{Seq}}
\newcommand{\SymSeq}{\mathrm{SymSeq}}
\newcommand{\bisymseq}{\mathrm{biSymSeq}}
\newcommand{\SymFun}{\mathrm{SymFun}}
\newcommand{\PSm}{\mathrm{PSm}}
\newcommand{\Sm}{\mathrm{Sm}}
\newcommand{\Psh}{\mathrm{Psh}}
\newcommand{\Ind}{\mathrm{Ind}}
\newcommand{\mP}{P_1}
\newcommand{\Tw}{\mathrm{Tw}}
\newcommand{\Adj}{\mathsf{Adj}}
\newcommand{\Adjffr}{\mathsf{Adj}^{\operatorname{R}}_{\mathrm{ff}}}
\newcommand{\Adjffl}{\mathsf{Adj}^{\operatorname{L}}_{\mathrm{ff}}}

\newcommand{\Arlax}{\mathrm{Ar}^\mathrm{lax}}
\newcommand{\Arr}{\mathsf{Ar}}
\newcommand{\Conf}{\mathsf{Conf}}
\newcommand{\Ret}{\mathsf{Ret}}
\newcommand{\BFib}{\mathsf{BFib}}
\newcommand{\soft}[1]{[\![#1]\!]}
\newcommand{\LM}{\mathrm{LM}}
\newcommand{\SegCoPsh}{\mathrm{co}\mathcal{PS}\mathrm{h}}
\newcommand{\LFib}{\mathcal{LF}\mathrm{ib}}
\newcommand{\Seg}{\mathcal{S}\mathrm{eg}}
\newcommand{\CSeg}{\mathcal{CS}\mathrm{eg}}
\newcommand{\un}{\smallint}
\newcommand{\st}{\widetilde}
\newcommand{\cart}{\mathrm{cart}}
\newcommand{\loc}{\mathrm{loc}}
\newcommand{\alln}{\Delta^\op_{/[n]}}
\newcommand{\ddelta}{{\mathlarger{\mathbbl{\Delta}}}}
\newcommand{\idln}{\Lambda^\op_{/[n]}}
\newcommand{\all}[1]{\Delta^\op_{/[#1]}}
\newcommand{\idl}[1]{\Lambda^\op_{/[#1]}}
\newcommand{\Corr}{\mathcal{C}\mathrm{orr}}
\newcommand{\twoCorr}{\mathrm{Corr}}
\newcommand{\CondFib}{\mathcal{C}\mathrm{ond}}
\newcommand{\Op}{\mathcal{O}\mathrm{p}}
\newcommand{\Fbrs}{\mathcal{F}\mathrm{brs}}
\newcommand{\WSF}{\mathcal{W}\mathrm{SF}}

%% Typical sub/superscripts
\mathchardef\mhyphen="2D
\newcommand{\cocart}{\mathrm{cocart}}
\newcommand{\even}{\mathrm{even}}
\newcommand{\odd}{\mathrm{odd}}
\newcommand{\cocomp}{\mathrm{cocomp}}
\newcommand{\comp}{\mathrm{comp}}
\newcommand{\lex}{\mathrm{lex}}
\newcommand{\ccoc}{\mhyphen\mathrm{coc}}
\newcommand{\coc}{\mathrm{coc}}
\newcommand{\bifib}{\mathrm{bifib}}
\newcommand{\lax}{\mathrm{lax}}
\newcommand{\oplax}{\mathrm{oplax}}
\newcommand{\sep}{\mathsf{sep}}
\newcommand{\cell}{\mathsf{cell}}
\newcommand{\intv}{\mathsf{int}}
\newcommand{\fintv}{\mathsf{fnt}}
\newcommand{\laxmon}{{\otimes\textrm{-}\mathrm{lax}}}
\newcommand{\oplaxmon}{\otimes\textrm{-}\mathrm{oplax}}



\newcommand{\ra}{\mathrm{ra}}
% Functors
\newcommand{\ver}{\mathrm{v}}
\newcommand{\hor}{\mathrm{h}}
\newcommand{\hop}{\mathrm{hop}}
\newcommand{\eL}{\mathcal{L}}
\newcommand{\eR}{\mathcal{R}}
\newcommand{\shiftcone}[1]{{#1}^\curlyvee}

%% bold font
\newcommand{\N}{\mathbf{N}}
\newcommand{\Z}{\mathbf{Z}}
\newcommand{\F}{\mathbf{F}}
\newcommand{\Q}{\mathbf{Q}}
\newcommand{\LL}{\mathbf{L}}
\newcommand{\Qbar}{\overline{\Q}}
\newcommand{\R}{\mathbf{R}}
\newcommand{\C}{\mathbf{C}}
\newcommand{\E}{\mathbf{E}}
\newcommand{\Sph}{\mathbb{S}}
\newcommand{\unit}{\mathds{1}}

%% mathcal font
\newcommand{\cA}{\mathcal{A}}
\newcommand{\cB}{\mathcal{B}}
\newcommand{\cC}{\mathcal{C}}
\newcommand{\cD}{\mathcal{D}}
\newcommand{\cE}{\mathcal{E}}
\newcommand{\cF}{\mathcal{F}}
\newcommand{\cG}{\mathcal{G}}
\newcommand{\cH}{\mathcal{H}}
\newcommand{\cI}{\mathcal{I}}
\newcommand{\cJ}{\mathcal{J}}
\newcommand{\cK}{\mathcal{K}}
\newcommand{\cL}{\mathcal{L}}
\newcommand{\cM}{\mathcal{M}}
\newcommand{\cN}{\mathcal{N}}
\newcommand{\cO}{\mathcal{O}}
\newcommand{\cP}{\mathcal{P}}
\newcommand{\cQ}{\mathcal{Q}}
\newcommand{\cR}{\mathcal{R}}
\newcommand{\cS}{\mathcal{S}}
\newcommand{\cT}{\mathcal{T}}
\newcommand{\cU}{\mathcal{U}}
\newcommand{\cV}{\mathcal{V}}
\newcommand{\cW}{\mathcal{W}}
\newcommand{\cX}{\mathcal{X}}
\newcommand{\cY}{\mathcal{Y}}
\newcommand{\cZ}{\mathcal{Z}}
\newcommand{\sX}{\mathcal{X}}
\newcommand{\sY}{\mathcal{Y}}
\newcommand{\sZ}{\mathcal{Z}}
\newcommand{\sU}{\mathcal{U}}
\newcommand{\sM}{\mathcal{M}}

% Various arrows
\makeatletter
\newcommand{\oset}[3][0ex]{%
	\mathrel{\mathop{#3}\limits^{
			\vbox to#1{\kern-2\ex@
				\hbox{$\scriptstyle#2$}\vss}}}}
\makeatother
\newcommand{\eqarrow}{\oset[-.16ex]{\sim}{\longrightarrow}}
\newcommand{\eqarrowback}{\oset[-.16ex]{\sim}{\longleftarrow}}
\tikzset{
	rot90/.style={anchor=south, rotate=90, inner sep=.5mm}
} % Rotating the \sim symbol in a tikz label bij 90 degrees.

% Overline enzo
\newcommand{\ol}{\overline}
\newcommand{\wt}{\widetilde}
\newcommand{\ul}{\underline}

% QED end of line displaymath theorem environment
\newcommand{\thmqed}[1]{\pushQED{\qed} #1 \qedhere\popQED}

% Subset symbols replaced by subseteq
\renewcommand{\subset}{\subseteq}
\renewcommand{\supset}{\supseteq}
\renewcommand{\epsilon}{\varepsilon}

\setlength{\parindent}{0.0in}
\setlength{\parskip}{0.05in}
% arrows and adjunctions
\def\lra{\longrightarrow}
\def\lla{\longleftarrow}
\def\hra{\hookrightarrow}
\def\llra{\def\arraystretch{.1}\begin{array}{c} \lra \\ \lla \end{array}}


\DeclareMathOperator{\li}{\mathsf{li}}



%%%%% arrowslashes

\makeatletter
\newcommand{\fixed@sra}{$\vrule height 2\fontdimen22\textfont2 width 0pt\rightarrow$}
\newcommand{\shortarrowup}[1]{%
  \mathrel{\text{\rotatebox[origin=c]{65}{\fixed@sra}}}
}
\newcommand{\shortarrowdown}[1]{%
  \mathrel{\text{\rotatebox[origin=c]{250}{\fixed@sra}}}
}
\makeatother

\newcommand{\upslash}{\!\shortarrowup{1}}
\newcommand{\upslashkern}{\!\!\!\!\shortarrowup{1}}
\newcommand{\downslash}{\!\shortarrowdown{1}}
